% ===============================================================
\chapter{Grice's Theory of Implicature}
% ===============================================================

\section{Paul Grice}

Grice was a defender of the following:
\begin{enumerate}[label=\roman*.]
  \item Natural language meaning is reducible to speakers' intentions, and thus to speakers' psychology (since intentions are mental states). He advocated for psychologism.
  \item Not all aspects of natural language meaning are conventional. There is also an aspect contributed by what speakers intend to convey in a given context, consistent with certain general rules governing conversation.
\end{enumerate}

\section{Grice's categorisation of Strawson, Frege and Russell}

\begin{definition}[Formalists]
Grice thought that Frege and Russell were \emph{formalists}. Formalists think that natural language arguments should be translated into a system of formal logic, where the validity or invalidity of an argument depends on whether statements of the same form generally preserve the truth of a conclusion given the truth of the premises. Formalists construct ideal logical languages.
\end{definition}

\begin{definition}[Informalists]
Grice thought that Strawson was an \emph{informalist}. Informalists think that the construction of an ideal logical language rests on certain faulty assumptions:
\begin{enumerate}[label=\roman*)]
  \item that the adequacy of language is to be determined by how well it can serve the purposes of science;
  \item that an expression is intelligible only if it can be regimented in a manner analysable in terms of symbols like `$\neg$' or `$\wedge$';
  \item that definitions are acceptable only if they can be expressed as logical equivalences.
\end{enumerate}
Many arguments in natural language are valid even if they do not fit a valid logical form, and so a less formal, more simplified way of determining acceptable forms of reasoning is required.
\end{definition}

\section{Implicature}

\begin{definition}[Implicature]
If $q$ is an \emph{implicature} of $p$, then $q$ could be either true or false when $p$ is true. Compare entailment: if $q$ is an entailment of $p$, then $q$ cannot be false when $p$ is true.
\end{definition}

\begin{itemize}
  \item Entailment is a \emph{semantic} relation. It concerns what is true of the conventional meaning of the sentence.
  \item Implicature is a \emph{pragmatic} relation. It concerns what we can infer based on a speaker's likely intentions given the circumstances of an utterance.
  \item Test: if $p \wedge q$ and $\neg q \wedge p$ can both occur, then q is an implicature of p. Else $p \wedge q$ and $\neg p \wedge \neg q$ and $\neg p \wedge q$ is an entailment.
  \item Implicatures are also \emph{cancelable}, meaning they can be taken back without a contradiction (i.e.\ $\neg q$ without implying $\neg p$).
\end{itemize}

\section{Types of implicature and Grice's maxims}

\begin{definition}[Types of implicature]\leavevmode
\begin{enumerate}[label=\roman*.]
  \item \emph{Conventional}: implicatures resulting from the conventional meanings of words or phrases, e.g.\ `but' implicates contrast.
  \item \emph{Conversational}: implicatures resulting from assumptions required to ensure that speakers are being cooperative (or why they aren't).
\end{enumerate}
\end{definition}

\begin{definition}[Cooperative principle]
The main maxim for conversational implicatures: make your conversational contribution such as is required, at the stage at which it occurs, by the accepted purpose or direction of the talk exchange in which you are engaged.
\end{definition}

Sub-maxims of the cooperative principle:
\begin{enumerate}[label=\roman*.]
  \item \emph{Maxim of quantity}: provide the right amount of information (not too little, not much more).
  \item \emph{Maxim of quality}: be accurate (try to make the statement true).
  \item \emph{Maxim of relation}: be relevant.
  \item \emph{Maxim of manner}: be perspicuous (be clear).
\end{enumerate}

\section{Meaning according to Grice}

\begin{definition}\leavevmode
\begin{itemize}
  \item \emph{Conventional meaning} (or \emph{non-natural meaning}) is what is conveyed by linguistic and cultural conventions like sarcasm etc.
  \item \emph{Speaker meaning} is the meaning conveyed by the utterance itself (conveyed meaning).
  \item \emph{Natural meaning}: meaning which is not conventional.
\end{itemize}
\end{definition}

Metalinguistic negation (``I didn't just giggle, I LAUGHED''), sarcasm and irony are examples where conventional meaning and speaker meaning differ.

\begin{definition}[Invited inferences]
Inferences hearers are invited to make given what was said and the circumstances in which an utterance takes place (asked to evaluate if there is an inference).
\begin{itemize}
  \item These are a subset of conversational implicatures.
  \item These have to be both calculable and cancelable.
\end{itemize}
\end{definition}

\begin{example}\leavevmode
\begin{enumerate}
  \item An utterance of `Mary and John fell in love and got married' invites the inference that they fell in love before getting married.
  \item An utterance of `Mary and John fell in love and got married' implicates that they fell in love before getting married. Again, the order in which conjoined clauses are uttered often suggests that one action described happened before another. Thus, `Mary and John got married and fell in love' suggests that they got married before falling in love. Thus, the fact that a speaker chose to order the clauses in the manner they did suggests they intend to convey that the events described happened in the order suggested, and the Maxim of Manner requires that speakers be orderly. Moreover, the implicature is cancelable: one can consistently follow this up with `Oh, and by the way, they got married before falling in love'.
\end{enumerate}
\end{example}

\begin{definition}[Calculability]
We must be able to give an argument based on Grice's maxims where obeying or failing to obey a certain maxim or maxims serves as a premise, and the implicature generated serves as a conclusion.
\begin{itemize}
  \item All calculable inferences are conversational, and vice versa.
\end{itemize}
\end{definition}

\section{Grice's Razor}

\begin{definition}[Grice's Razor]
Do not postulate ambiguities beyond necessity.
\end{definition}

In other words, if we can explain multiple interpretations of a sentence pragmatically (via implicature) rather than semantically (by postulating an ambiguity in any of the words contained), then we should do that.

Third explanation: if there is an ambiguity, use an implicature to explain it.
\begin{itemize}
  \item Furthermore, this puts Strawson's theory of meaning in jeopardy, because Grice suggests that sentences without ``any truth value'' can be designated true or false regardless, thanks to implicature. For example, ``The present king of France is bald'' can be followed up with ``Wait a minute, there isn't a king of France.''
\end{itemize}

\section{Objections to Grice's theory}

According to Grice, sentences have meaning due to speakers uttering them with an intention to spark a corresponding belief in their audience. However, this collapses in the following scenarios:
\begin{enumerate}[label=\roman*.]
  \item \emph{Audienceless utterances}: talking to myself.
  \item \emph{Absence of intention}: saying something without the intent of inducing a belief, such as a student answering a question for somebody who already knows the answer and is convinced of the belief.
  \item \emph{Opposite intention}: saying things to inspire a belief contrary to what was said.
  \item \emph{Recursion}: there are indefinitely many meaningful sentences in natural language, and it is impossible to utter them all with some intention.
\end{enumerate}
